\chapter{Podsumowanie i kierunki dalszego rozwoju}
\label{chap:podsumowanie}

Przedmiotem pracy było zaprojektowanie i implementacja systemu PlanWise, wspierającego organizację pracy projektowej oraz przygotowywanie informacji pomocnych w podejmowaniu decyzji. Rozwiązanie łączy zarządzanie projektami, zespołami, zadaniami i sprintami z analizą ryzyka, priorytetyzacją backlogu, harmonogramowaniem, proponowaniem przydziałów oraz estymacją kosztów.

W niniejszym rozdziale podsumowano zakres opracowanego rozwiązania, wskazano wkład projektowy i implementacyjny oraz przedstawiono ograniczenia i kierunki rozwoju. Ocena odnosi się do funkcji obecnych w analizowanej wersji aplikacji. Wnioski dotyczące skuteczności rekomendacji wymagają wyników badań opisanych w poprzednim rozdziale.

\section{Stopień realizacji celu i zakresu pracy}

Cel projektowy zrealizowano przez opracowanie aplikacji internetowej z backendem opartym na .NET i frontendem utworzonym w Angularze. Część serwerową zorganizowano jako monolit modularny z podziałem odpowiedzialności pomiędzy osiem modułów.

Zaimplementowano podstawowe funkcje organizacji pracy: obsługę kont, tworzenie projektów, zarządzanie zespołem, definiowanie zadań, prowadzenie backlogu i obsługę sprintów. Tablica umożliwia wizualizację statusów i przenoszenie zadań, natomiast harmonogram przedstawia planowane terminy i zależności.

Dane operacyjne wykorzystano jako wspólne źródło informacji dla mechanizmów wspomagania decyzji. Dzięki temu użytkownik nie musi przygotowywać odrębnego zestawu danych dla każdej analizy. Jednocześnie jakość rekomendacji zależy od kompletności informacji wprowadzanych podczas pracy nad projektem.

Zakres realizacji poszczególnych obszarów można podsumować następująco:

\begin{enumerate}
    \item Organizacja i monitorowanie pracy obejmują projekty, zadania, sprinty, tablicę, harmonogram oraz wybrane zestawienia postępu. Nie wszystkie elementy pełnego zarządzania przepływem zostały zaimplementowane; przykładem są konfigurowalne limity WIP.

    \item Analiza ryzyka jest dostępna w dwóch wariantach: deterministycznej punktacji oraz modelu regresji logistycznej dopasowanego do zewnętrznego zbioru projektów zwinnych. Obie metody działają za wspólną granicą programistyczną i pozostają jednocześnie uruchamialne. System rejestruje ponadto wektory cech i wyniki realizacji, co docelowo umożliwi rekalibrację modelu na danych wdrożenia. Prognoza sprintu pozostaje projekcją opartą na tempie zespołu i nie pochodzi z modelu uczonego.

    \item Priorytetyzacja backlogu wykorzystuje ważoną ocenę wartości biznesowej, zależności, złożoności i ryzyka, a alternatywnie porządkowanie z udziałem modelu językowego. Wynik zawiera proponowaną kolejność i uzasadnienia.

    \item Harmonogramowanie obejmuje obliczenia oparte na zależnościach, a przydział zadań realizuje solver programowania z ograniczeniami, minimalizujący długość projektu przy zachowaniu kolejności zależności i rozłączności prac jednej osoby. Nie uwzględniono natomiast kalendarzy nieobecności, a jeden punkt nadal odpowiada jednemu dniu pracy przy pełnej dostępności.

    \item Estymacja kosztów korzysta z rzeczywistej integracji z modelem językowym i danych zespołu. Przetwarzanie odpowiedzi wymaga jednak rozbudowy kontroli semantycznej i rachunkowej, w tym zapewnienia kompletności wymaganych wariantów kosztowych.
\end{enumerate}

Zrealizowano zatem aplikację integrującą funkcje operacyjne i analityczne, przy czym część pierwotnych założeń pozostaje spełniona w ograniczonym zakresie. Nie dotyczy to już samego istnienia modelu uczonego ani optymalizacji uwzględniającej ograniczenia zasobowe, lecz wykazania ich przewagi w warunkach konkretnego wdrożenia oraz pełnego odwzorowania kalendarza pracy.

\section{Najważniejsze rezultaty projektowe i implementacyjne}

Rezultatem pracy jest spójny model aplikacji, w którym dane projektu są wykorzystywane przez kilka mechanizmów wspomagania decyzji. Zadanie stanowi jednocześnie element backlogu, składnik harmonogramu i źródło cech dla analiz ryzyka, priorytetów oraz kosztów.

Podział na moduły pozwolił oddzielić odpowiedzialność za dane operacyjne od odpowiedzialności za ich analizowanie. Wspólne kontrakty umożliwiają przekazywanie potrzebnych informacji bez bezpośredniego korzystania z wewnętrznych encji innych modułów.

Opracowano również wspólny sposób uruchamiania analiz w tle. Zlecenia posiadają identyfikatory i stan wykonania, a klient może pobrać wynik po zakończeniu operacji. Takie rozwiązanie obsługuje zarówno lokalne obliczenia, jak i oczekiwanie na odpowiedź modelu językowego.

Istotnym elementem interfejsu jest rozdzielenie propozycji od jej zastosowania. Użytkownik może przejrzeć rekomendację i uzasadnienie przed wykonaniem zmiany. W przypadku estymacji kosztów możliwe jest również odwołanie do wcześniejszych wyników.

Zastosowanie jawnych funkcji heurystycznych umożliwia prześledzenie sposobu wyznaczania ocen. Ma to znaczenie podczas weryfikacji obliczeń i projektowania eksperymentów. Ujawnia także zależność wyniku od przyjętych wag, progów i wartości zastępczych.

Odrębnym rezultatem jest ujednolicenie sposobu, w jaki system traktuje metody wyznaczania wyników. Wszystkie cztery mechanizmy analityczne umieszczono za granicami programistycznymi o tym samym kształcie, a za każdą z nich zarejestrowano co najmniej dwie implementacje. Wraz z każdym uruchomieniem zapisywana jest nazwa metody, która je wykonała, oraz przyjęte przez nią założenia.

Konstrukcja ta nie jest wyłącznie zabiegiem porządkującym kod. Umożliwia porównywanie metod zamiast zakładania przewagi którejkolwiek z nich, a w przypadku harmonogramowania pozwala obsłużyć niepowodzenie solvera jako przewidziany przypadek, którego skutkiem jest wynik gorszy, lecz jawnie oznaczony, a nie błąd analizy.

Dodatkowym rezultatem jest uruchomienie rejestrowania danych uczących. Element ten ma charakter nieodwracalny w czasie: stan projektu z danego dnia nie może zostać odtworzony później, więc rozpoczęcie zapisu jest warunkiem, a nie następstwem dalszych prac nad predykcją opartą na danych własnych.

W obszarze modelowania rezultatem jest wytrenowany i wdrożony model predykcji opóźnienia zadania wraz z udokumentowaną procedurą jego przygotowania, obejmującą dobór zbioru danych, definicję etykiety, usunięcie wykrytych źródeł przecieku informacji, wybór postaci cech oraz kalibrację. Wynik pomiaru przedstawiono bez upiększeń: przewaga nad metodą odniesienia jest rzeczywista, lecz umiarkowana, co bezpośrednio uzasadniło pozostawienie heurystyki jako metody domyślnej.

Powyższe rezultaty dotyczą opracowanych mechanizmów i ich organizacji. Nie stanowią jeszcze potwierdzenia zmniejszenia opóźnień, kosztów lub czasu pracy kierownika projektu.

\section{Wkład własny}

Wkład własny obejmuje zaprojektowanie i połączenie funkcji zarządzania projektami z mechanizmami analitycznymi w ramach jednej aplikacji. Dotyczy zarówno modelu danych i granic modułów, jak i realizacji przypadków użycia oraz sposobu prezentacji rekomendacji.

W części serwerowej wkład obejmuje implementację obsługi projektów, zespołów i zadań, integrację modułów przez wspólne kontrakty oraz przygotowanie przebiegu analiz asynchronicznych. W części klienckiej obejmuje widoki organizacji pracy i obsługę wyników analiz w kontekście wybranego projektu.

W obszarze analitycznym opracowano konfigurację czynników ryzyka, funkcję priorytetyzacji i heurystykę wyboru wykonawców. Przygotowano również kontekst estymacji kosztów oparty na danych backlogu i zagregowanych stawkach zespołu oraz integrację odpowiedzi modelu z aplikacją.

W zakresie prac nad modelami wkład obejmuje zaprojektowanie jednolitej granicy oddzielającej sposób wyznaczania wyniku od pozostałych odpowiedzialności modułu i zastosowanie jej we wszystkich czterech mechanizmach, sformułowanie zdarzenia docelowego i sposobu jego etykietowania, dobór rejestrowanych cech oraz opracowanie mechanizmu gromadzenia i etykietowania obserwacji wraz z obsługą przypadków cenzurowanych.

Wkład w części dotyczącej modelu uczonego obejmuje ponadto dobór zbioru danych wraz z rozpoznaniem przyczyn nieprzydatności zbioru odrzuconego, przedefiniowanie zobowiązania terminowego na datę zakończenia sprintu, wykrycie i usunięcie dwóch źródeł przecieku informacji, zaprojektowanie dwóch protokołów oceny, przeprowadzenie porównania rodzin modeli i postaci cech, ograniczenie zestawu cech do wielkości możliwych do obliczenia podczas eksploatacji oraz kalibrację wyniku. Osobnym elementem jest zabezpieczenie zgodności dwóch implementacji tych samych obliczeń zestawem przypadków wzorcowych.

W obszarze harmonogramowania wkład dotyczy sformułowania przydziału zadań jako problemu harmonogramowania z ograniczonymi zasobami, doboru ograniczeń i postaci funkcji celu oraz zaprojektowania sposobu obsługi przypadków, w których solver nie zwraca rozwiązania.

W odniesieniu do modeli językowych wkład dotyczy sformułowania zadania, ograniczenia odpowiedzi do ról i stawek pochodzących z danych projektu, wymuszenia ustrukturyzowanej postaci wyniku, rozróżnienia niepowodzeń wymagających ponowienia od tych, dla których jest ono bezcelowe, oraz rozpoznania, że priorytetyzacja backlogu jest problemem pozbawionym etykiety, a przez to właściwym raczej dla modelu językowego niż dla modelu uczonego.

Wykorzystane podstawy, takie jak metoda ścieżki krytycznej, suma ważona czy podejście zachłanne, są znanymi metodami. Wkład pracy polega na ich doborze, dostosowaniu i integracji z modelem PlanWise, a nie na opracowaniu nowej klasy algorytmów.

Elementem opracowania jest także wskazanie ograniczeń interpretacji wyników. Rozróżnienie punktacji od prawdopodobieństwa, scenariusza od kwantyla i deklarowanej oszczędności od rzeczywistego efektu pozwala precyzyjniej określić możliwości rozwiązania.

\section{Ocena tezy pracy}

Teza pracy dotyczy możliwości poprawy wspomagania decyzji przez połączenie danych projektowych z mechanizmami analitycznymi. Do jej oceny potrzebne jest porównanie PlanWise z określonymi metodami bazowymi w tych samych warunkach.

Opracowana implementacja pokazuje techniczną możliwość przeprowadzenia takiej integracji. System przygotowuje oceny i propozycje na podstawie danych wykorzystywanych podczas codziennej organizacji pracy, a użytkownik może odnieść wynik do konkretnych zadań i członków zespołu.

Jeden element tezy uzyskał wsparcie empiryczne. Dla predykcji opóźnień wykazano pomiarem przewagę modelu uczonego nad metodą odniesienia w protokole odpowiadającym warunkom rzeczywistego użycia, to znaczy na projektach nieobecnych w zbiorze uczącym. Przewaga ta jest jednak umiarkowana, a uzyskane wartości uprawniają do porządkowania uwagi kierownika, nie zaś do orzekania o losie pojedynczego zadania.

Wyniku tego nie należy przy tym utożsamiać z weryfikacją tezy w jej pełnym brzmieniu. Metodą odniesienia był model stały, a nie heurystyka stosowana w systemie, wobec czego wykazano, że w danych istnieje sygnał możliwy do wykorzystania, a nie że metoda uczona przewyższa dostrojoną regułę ekspercką. Porównanie tych dwóch metod pozostaje otwarte i wymaga danych gromadzonych przez wdrożenie.

Dla pozostałych mechanizmów brakuje wyników dotyczących użyteczności rankingu, jakości przydziałów i błędów estymacji kosztów. Jeżeli teza obejmuje również usprawnienie pracy użytkownika, konieczne jest zbadanie tego efektu w zadaniach decyzyjnych.

Teza nie może zatem zostać uznana za potwierdzoną w całości. Nie została też odrzucona: uzyskano cząstkowy wynik pozytywny oraz wykazano wykonalność integracji, natomiast zakres możliwego wnioskowania pozostaje ograniczony brakiem pozostałych pomiarów.

% DO AKTUALIZACJI PO UZUPEŁNIENIU ROZDZIAŁU 7:
% Wskazać, które wyniki wspierają tezę, w jakich warunkach
% i dla jakich mechanizmów. Uwzględnić również wyniki negatywne.
% Nie utożsamiać poprawności implementacji z jakością decyzji.

\section{Ograniczenia opracowanego rozwiązania}

Najważniejszym ograniczeniem pozostaje pochodzenie danych uczących. Model dopasowano do projektów otwartego oprogramowania pochodzących z innych organizacji, które mogą opóźniać się inaczej niż zespoły komercyjne działające pod presją zobowiązań wobec klienta. Prawdopodobieństwa nie zostały skalibrowane na wynikach tego wdrożenia, wobec czego należy je odczytywać jako sposób porządkowania uwagi, a nie jako dosłowne szanse wystąpienia zdarzenia.

Mechanizm gromadzenia danych własnych został zaimplementowany i działa przy każdym uruchomieniu analizy ryzyka, jednak zbiór narasta dopiero wraz z eksploatacją systemu. Ograniczenia tego nie da się usunąć zmianą w kodzie: stan wiedzy dostępny w chwili wcześniejszych decyzji nie może zostać odtworzony wstecz, a jedynie zapisany na bieżąco. W konsekwencji przewaga modelu uczonego nad heurystyką stosowaną w systemie nie została wykazana i to właśnie ta okoliczność, a nie względy techniczne, decyduje o pozostawieniu heurystyki jako metody domyślnej.

Ograniczony jest również zakres tego, co model przewiduje. Wyznacza on wyłącznie prawdopodobieństwo opóźnienia zadania, natomiast wielkość opóźnienia oraz prognoza wykonania sprintu pozostają wynikiem przeliczeń heurystycznych, ponieważ wykorzystany zbiór nie zawiera odpowiednich etykiet.

Drugim ograniczeniem jest sposób reprezentowania czasu i zasobów. Punkty, godziny, dni kalendarzowe i udział etatu opisują różne wielkości, a ich powiązanie wymaga jawnego przelicznika. Zagadnienie to rozwiązano dla prognozowania sprintu, wprowadzając wyznaczanie tempa z zakończonych sprintów projektu wraz z jawnym rozróżnieniem tempa zmierzonego od założonego. Nie rozwiązano go natomiast dla czasu trwania zadania, gdzie utrzymano uproszczenie jednego punktu jako jednego dnia pracy przy pełnej dostępności.

Analiza zależności wymaga również pełnej kontroli acykliczności. Przygotowanie danych wykresu mimo występowania cyklu nie jest równoznaczne z wyznaczeniem wykonalnego harmonogramu.

Przydział zadań uwzględnia obecnie kolejność zależności, rozłączność prac jednej osoby i skalowanie czasu trwania dostępnością, nie modeluje natomiast kalendarzy nieobecności ani zmienności dostępności w czasie. Dopasowanie kompetencji pozostaje porównaniem tekstowym z tytułem zadania i nie uwzględnia poziomu umiejętności, przez co pełni w funkcji celu wyłącznie rolę rozstrzygania remisów. Optymalność jest gwarantowana jedynie względem zadeklarowanej funkcji celu i tylko wtedy, gdy solver zdąży jej dowieść w wyznaczonym czasie.

Estymacja kosztów jest ograniczona zakresem przekazywanego kontekstu i brakiem danych o rzeczywistym nakładzie pozostałym do wykonania. Poprawność struktury odpowiedzi modelu nie gwarantuje poprawności sum, kompletności scenariuszy ani trafności uzasadnienia.

Dodatkowe ograniczenia dotyczą aktualności rekomendacji i niezawodności ich stosowania. Dane mogą zmienić się między wygenerowaniem propozycji a jej zatwierdzeniem. Operacje obejmujące kilka modułów nie są automatycznie jedną transakcją, a obsługa wielu instancji procesorów wymaga dodatkowych mechanizmów współbieżności.

Ograniczenia te wyznaczają zakres zastosowania obecnej wersji. System należy traktować jako rozwiązanie rozwojowe, którego rekomendacje wymagają oceny użytkownika i dalszej walidacji.

\section{Kierunki dalszego rozwoju}

\subsection{Ujednolicenie danych i walidacji}

Pierwszym kierunkiem rozwoju powinno być uporządkowanie modelu pracochłonności i dostępności. Wskazane jest rozdzielenie oszacowania punktowego, planowanych godzin, pozostałego nakładu i czasu kalendarzowego.

Dostępność zespołu powinna odnosić się do określonego okresu i uwzględniać dni robocze oraz nieobecności. Pozwoli to porównywać obciążenie i pojemność w tych samych jednostkach.

Należy również wprowadzić kontrolę cykli zależności, pełniejszą walidację wyników kosztowych i oznaczanie brakujących danych. Nazwy pól sugerujące prawdopodobieństwa lub kwantyle powinny odpowiadać rzeczywistemu charakterowi obliczeń.

\subsection{Gromadzenie historii i rozwój predykcji}

Model predykcji opóźnień został wytrenowany i wdrożony, dlatego dalsze prace dotyczą już nie jego zbudowania, lecz uzasadnienia jego użycia w konkretnym wdrożeniu. Krokiem rozstrzygającym jest rekalibracja na danych gromadzonych przez system oraz porównanie z heurystyką na tych samych obserwacjach testowych, z podziałem uwzględniającym zarówno porządek czasowy, jak i przynależność obserwacji do zadań. Dopiero wynik takiego porównania powinien decydować o zmianie metody domyślnej.

Poprzedzać go musi scharakteryzowanie zgromadzonego zbioru, obejmujące jego liczebność, rozkład etykiet oraz udział obserwacji cenzurowanych. Wielkości tych nie da się przewidzieć przed zebraniem danych, a mogą one wymusić zmianę doboru metryk lub sposobu uczenia.

Rozszerzenia wymaga również sam zakres rejestrowanych danych. Obecny zapis obejmuje stan zadania w chwili prognozy, nie zawiera natomiast historii zmian terminów, oszacowań i przypisań pomiędzy kolejnymi uruchomieniami analizy. Uzupełnienie go pozwoliłoby rozpoznać zadania, których parametry były wielokrotnie korygowane, co samo w sobie może stanowić sygnał zagrożenia. Odrębną luką jest brak historii przepustowości zespołu, w której prawdopodobnie znajduje się istotna część sygnału nieuchwyconego przez obecny zestaw cech.

Kalibrację zastosowano dla modelu uczonego, natomiast nie objęto nią heurystyki, której wynik zapisywany jest w polach o nazwach odwołujących się do prawdopodobieństwa. Uporządkowanie tego nazewnictwa lub przeprowadzenie kalibracji pozostaje do wykonania. Niezależnie od wybranej metody konieczne jest monitorowanie jakości prognoz w czasie, ponieważ w danych zaobserwowano rzeczywistą zmianę udziału opóźnień pomiędzy okresami.

\subsection{Rozbudowa harmonogramowania i przydziałów}

Zadanie optymalizacyjne z ograniczeniami zasobowymi i czasowymi zostało sformułowane i wdrożone, a heurystyka zachłanna pozostała szybką metodą bazową i rozwiązaniem zastępczym. Dalszy rozwój dotyczy zatem przede wszystkim wierności modelu wobec rzeczywistości.

Pierwszym kierunkiem jest strukturalny opis wymagań kompetencyjnych zadań, pozwalający zastąpić dopasowanie fragmentów tytułu jawnymi warunkami wyboru wykonawcy. Obecna postać tego składnika jest na tyle zawodna, że w funkcji celu przypisano mu wyłącznie rolę rozstrzygania remisów; wprowadzenie opisu strukturalnego pozwoliłoby podnieść jego znaczenie.

Drugim jest wprowadzenie kalendarzy członków zespołu, obejmujących nieobecności i zmienną dostępność w czasie. Obecnie każdy dzień traktowany jest jako roboczy, a dostępność jako wielkość stała, co przy dłuższym horyzoncie planowania prowadzi do terminów zbyt optymistycznych.

Trzecim jest zmierzenie praktycznej granicy stosowalności solvera. Należy ustalić, przy jakiej wielkości projektu rośnie udział uruchomień kończących się wynikiem metody zastępczej, a następnie rozważyć zwiększenie budżetu czasu, dekompozycję zadania lub rozwiązanie przyrostowe.

Rozważenia wymaga również sposób łączenia składników funkcji celu. Obecne wagi stałe dopuszczają teoretycznie wymianę dnia długości projektu na dostatecznie dużą oszczędność na przekroczeniach terminów; rozwiązanie leksykograficzne usunęłoby tę możliwość kosztem większej złożoności obliczeń.

\subsection{Zwiększenie wiarygodności estymacji kosztów}

Model językowy powinien koncentrować się na interpretacji wymagań, proponowaniu założeń i oszacowaniu nakładu. Obliczenia rachunkowe oraz kontrola zgodności stawek powinny być realizowane deterministycznie po stronie aplikacji.

Kontrakt wyniku warto rozszerzyć o jednoznaczne powiązanie pozycji kosztowych ze scenariuszem. Należy też sprawdzać zależności między rekomendacjami redukcji, aby uniknąć wielokrotnego uwzględniania tej samej oszczędności.

Dalszy rozwój obejmuje rejestrowanie dokładnego kontekstu wywołania, wersji instrukcji i konfiguracji modelu. Dane te powinny również wpływać na decyzję o ponownym wykorzystaniu wcześniejszego wyniku.

Po zgromadzeniu rzeczywistych kosztów możliwe będzie porównywanie estymacji z wykonaniem i identyfikowanie obszarów systematycznego niedoszacowania.

\subsection{Niezawodność i kontrola zmian}

Przed zastosowaniem propozycji należy ponownie sprawdzać jej aktualność. Można wykorzystać wersjonowanie danych lub zapis identyfikatora stanu, dla którego wykonano analizę.

Obsługa przydziałów powinna jednoznacznie raportować częściowe niepowodzenia. Procesory zadań i wiadomości wymagają natomiast zasad przejmowania pracy, ponowień i odporności na powtórne wykonanie.

Rozwój kontroli dostępu powinien obejmować jawne uprawnienia do poszczególnych działań oraz ich weryfikację w API i komunikacji czasu rzeczywistego.

\subsection{Badania z udziałem użytkowników}

Po usunięciu podstawowych ograniczeń wskazane jest przeprowadzenie badań z udziałem osób zarządzających projektami. Uczestnicy mogliby wykonywać te same zadania decyzyjne z dostępem do rekomendacji i bez nich.

Ocena powinna obejmować czas analizy, liczbę przeoczonych problemów, poprawność decyzji i zrozumienie uzasadnień. Należy również sprawdzić, czy sposób prezentacji nie prowadzi do nadmiernego zaufania do wyników.

Takie badanie pozwoli uzupełnić pomiary techniczne o ocenę rzeczywistej użyteczności. Dopiero połączenie tych perspektyw umożliwi określenie, w jakich warunkach PlanWise wspiera kierownika projektu i które mechanizmy wymagają dalszych zmian.